\section{Matrix factorizations}
An $LU$ factorization writes $A=LU$, where $L$ is unit lower
triangular and $U$ is in echelon form. ``Unit'' means the diagonal
entries of $L$ are $1$. In elimination, if we replace row $i$ by
$R_i-mR_j$ to clear an entry below the current pivot, record $m$
in the corresponding lower column of $L$. The same row replacements
change $A$ to $U$ and $L$ to $I$; reversing them proves $A=LU$.
No row interchanges are needed for the matrices factored below.

To solve $Ax=b$, first solve $Ly=b$ from top to bottom, and then
$Ux=y$ from bottom to top. Reuse $L,U$ when the right side changes;
do not compute an inverse merely to solve one system.

\practice{2.5.P1}{A rectangular matrix with only three pivot columns}
The given matrix has five rows and four columns. Its first-column
multipliers are $3,1,2,-3$. After that elimination, the second-column
multipliers are $-1,2,-3$. Column 3 then has no new pivot; the next
pivot is the entry $5$ in column 4 of row 3. Its multipliers are $-1,2$.
The resulting factors are

\[L=\left[\begin{matrix}1 & 0 & 0 & 0 & 0\\3 & 1 & 0 & 0 & 0\\1 & -1 & 1 & 0 & 0\\2 & 2 & -1 & 1 & 0\\-3 & -3 & 2 & 0 & 1\end{matrix}\right].\]
\[U=\left[\begin{matrix}2 & -4 & -2 & 3\\0 & 3 & 1 & -1\\0 & 0 & 0 & 5\\0 & 0 & 0 & 0\\0 & 0 & 0 & 0\end{matrix}\right].\]

The unused fourth and fifth columns of $L$ are those of $I_5$.
Multiplying these displayed factors gives the original $5\times4$
matrix exactly. In particular, $L$ is $5\times5$, not $5\times4$;
$U$ carries the rectangular dimensions and its last two rows are zero.

\exercise{2.5.1}{Forward substitution followed by backward substitution}
\goal Solve $Ax=b$ using the given factors.
\[A=\left[\begin{matrix}3 & -7 & -2\\-3 & 5 & 1\\6 & -4 & 0\end{matrix}\right],\qquad b=\left[\begin{matrix}-7\\5\\2\end{matrix}\right].\]
The supplied factors are
\[L=\left[\begin{matrix}1 & 0 & 0\\-1 & 1 & 0\\2 & -5 & 1\end{matrix}\right].\]
\[U=\left[\begin{matrix}3 & -7 & -2\\0 & -2 & -1\\0 & 0 & -1\end{matrix}\right].\]
\step{First solve $Ly=b$ by forward substitution.}
\begin{align*}y_{1}&=-7\\
y_{2}&=5-\bigl[(-1)(-7)\bigr]=-2\\
y_{3}&=2-\bigl[(2)(-7)+(-5)(-2)\bigr]=6\end{align*}\step{Then solve $Ux=y$ by backward substitution.}
\begin{align*}x_{3}&=\frac{6}{-1}=-6\\
x_{2}&=\frac{-2-\bigl[(-1)(-6)\bigr]}{-2}=4\\
x_{1}&=\frac{-7-\bigl[(-7)(4)+(-2)(-6)\bigr]}{3}=3\end{align*}\step{Also carry out ordinary elimination.}
Apply these row replacements to the augmented matrix $[A\mid b]$:
\begin{align*}R_{2}&\leftarrow R_{2}+R_{1}\\
R_{3}&\leftarrow R_{3}-(2)R_{1}\\
R_{3}&\leftarrow R_{3}+(5)R_{2}\end{align*}They produce $[U\mid y]$, with the same $U,y$ displayed above.
The backward substitution therefore gives the same vector as the
factorization method. This explicitly performs the requested second method.
\answer{$x=\left[\begin{matrix}3\\4\\-6\end{matrix}\right]$. Multiplying by the original $A$
returns the displayed $b$, an independent check of the solution.}

\exercise{2.5.2}{Forward substitution followed by backward substitution}
\goal Solve $Ax=b$ using the given factors.
\[A=\left[\begin{matrix}2 & -6 & 4\\-4 & 8 & 0\\0 & -4 & 6\end{matrix}\right],\qquad b=\left[\begin{matrix}2\\-4\\6\end{matrix}\right].\]
The supplied factors are
\[L=\left[\begin{matrix}1 & 0 & 0\\-2 & 1 & 0\\0 & 1 & 1\end{matrix}\right].\]
\[U=\left[\begin{matrix}2 & -6 & 4\\0 & -4 & 8\\0 & 0 & -2\end{matrix}\right].\]
\step{First solve $Ly=b$ by forward substitution.}
\begin{align*}y_{1}&=2\\
y_{2}&=-4-\bigl[(-2)(2)\bigr]=0\\
y_{3}&=6-\bigl[(1)(0)\bigr]=6\end{align*}\step{Then solve $Ux=y$ by backward substitution.}
\begin{align*}x_{3}&=\frac{6}{-2}=-3\\
x_{2}&=\frac{0-\bigl[(8)(-3)\bigr]}{-4}=-6\\
x_{1}&=\frac{2-\bigl[(-6)(-6)+(4)(-3)\bigr]}{2}=-11\end{align*}\step{Also carry out ordinary elimination.}
Apply these row replacements to the augmented matrix $[A\mid b]$:
\begin{align*}R_{2}&\leftarrow R_{2}+(2)R_{1}\\
R_{3}&\leftarrow R_{3}-R_{2}\end{align*}They produce $[U\mid y]$, with the same $U,y$ displayed above.
The backward substitution therefore gives the same vector as the
factorization method. This explicitly performs the requested second method.
\answer{$x=\left[\begin{matrix}-11\\-6\\-3\end{matrix}\right]$. Multiplying by the original $A$
returns the displayed $b$, an independent check of the solution.}

\exercise{2.5.3}{Forward substitution followed by backward substitution}
\goal Solve $Ax=b$ using the given factors.
\[A=\left[\begin{matrix}2 & -4 & 2\\-4 & 5 & 2\\6 & -9 & 1\end{matrix}\right],\qquad b=\left[\begin{matrix}6\\0\\6\end{matrix}\right].\]
The supplied factors are
\[L=\left[\begin{matrix}1 & 0 & 0\\-2 & 1 & 0\\3 & -1 & 1\end{matrix}\right].\]
\[U=\left[\begin{matrix}2 & -4 & 2\\0 & -3 & 6\\0 & 0 & 1\end{matrix}\right].\]
\step{First solve $Ly=b$ by forward substitution.}
\begin{align*}y_{1}&=6\\
y_{2}&=0-\bigl[(-2)(6)\bigr]=12\\
y_{3}&=6-\bigl[(3)(6)+(-1)(12)\bigr]=0\end{align*}\step{Then solve $Ux=y$ by backward substitution.}
\begin{align*}x_{3}&=0=0\\
x_{2}&=\frac{12-\bigl[(6)(0)\bigr]}{-3}=-4\\
x_{1}&=\frac{6-\bigl[(-4)(-4)+(2)(0)\bigr]}{2}=-5\end{align*}\answer{$x=\left[\begin{matrix}-5\\-4\\0\end{matrix}\right]$. Multiplying by the original $A$
returns the displayed $b$, an independent check of the solution.}

\exercise{2.5.4}{Forward substitution followed by backward substitution}
\goal Solve $Ax=b$ using the given factors.
\[A=\left[\begin{matrix}1 & -1 & 2\\1 & -3 & 1\\3 & 7 & 5\end{matrix}\right],\qquad b=\left[\begin{matrix}0\\-5\\7\end{matrix}\right].\]
The supplied factors are
\[L=\left[\begin{matrix}1 & 0 & 0\\1 & 1 & 0\\3 & -5 & 1\end{matrix}\right].\]
\[U=\left[\begin{matrix}1 & -1 & 2\\0 & -2 & -1\\0 & 0 & -6\end{matrix}\right].\]
\step{First solve $Ly=b$ by forward substitution.}
\begin{align*}y_{1}&=0\\
y_{2}&=-5-\bigl[(1)(0)\bigr]=-5\\
y_{3}&=7-\bigl[(3)(0)+(-5)(-5)\bigr]=-18\end{align*}\step{Then solve $Ux=y$ by backward substitution.}
\begin{align*}x_{3}&=\frac{-18}{-6}=3\\
x_{2}&=\frac{-5-\bigl[(-1)(3)\bigr]}{-2}=1\\
x_{1}&=0-\bigl[(-1)(1)+(2)(3)\bigr]=-5\end{align*}\answer{$x=\left[\begin{matrix}-5\\1\\3\end{matrix}\right]$. Multiplying by the original $A$
returns the displayed $b$, an independent check of the solution.}

\exercise{2.5.5}{Forward substitution followed by backward substitution}
\goal Solve $Ax=b$ using the given factors.
\[A=\left[\begin{matrix}1 & -2 & -2 & -3\\3 & -9 & 0 & -9\\-1 & 2 & 4 & 7\\-3 & -6 & 26 & 2\end{matrix}\right],\qquad b=\left[\begin{matrix}1\\6\\0\\3\end{matrix}\right].\]
The supplied factors are
\[L=\left[\begin{matrix}1 & 0 & 0 & 0\\3 & 1 & 0 & 0\\-1 & 0 & 1 & 0\\-3 & 4 & -2 & 1\end{matrix}\right].\]
\[U=\left[\begin{matrix}1 & -2 & -2 & -3\\0 & -3 & 6 & 0\\0 & 0 & 2 & 4\\0 & 0 & 0 & 1\end{matrix}\right].\]
\step{First solve $Ly=b$ by forward substitution.}
\begin{align*}y_{1}&=1\\
y_{2}&=6-\bigl[(3)(1)\bigr]=3\\
y_{3}&=0-\bigl[(-1)(1)\bigr]=1\\
y_{4}&=3-\bigl[(-3)(1)+(4)(3)+(-2)(1)\bigr]=-4\end{align*}\step{Then solve $Ux=y$ by backward substitution.}
\begin{align*}x_{4}&=-4=-4\\
x_{3}&=\frac{1-\bigl[(4)(-4)\bigr]}{2}=\frac{17}{2}\\
x_{2}&=\frac{3-\bigl[(6)(\frac{17}{2})\bigr]}{-3}=16\\
x_{1}&=1-\bigl[(-2)(16)+(-2)(\frac{17}{2})+(-3)(-4)\bigr]=38\end{align*}\answer{$x=\left[\begin{matrix}38\\16\\\frac{17}{2}\\-4\end{matrix}\right]$. Multiplying by the original $A$
returns the displayed $b$, an independent check of the solution.}

\exercise{2.5.6}{Forward substitution followed by backward substitution}
\goal Solve $Ax=b$ using the given factors.
\[A=\left[\begin{matrix}1 & 3 & 2 & 0\\-2 & -3 & -4 & 12\\3 & 0 & 4 & -36\\-5 & -3 & -8 & 49\end{matrix}\right],\qquad b=\left[\begin{matrix}1\\-2\\-1\\2\end{matrix}\right].\]
The supplied factors are
\[L=\left[\begin{matrix}1 & 0 & 0 & 0\\-2 & 1 & 0 & 0\\3 & -3 & 1 & 0\\-5 & 4 & -1 & 1\end{matrix}\right].\]
\[U=\left[\begin{matrix}1 & 3 & 2 & 0\\0 & 3 & 0 & 12\\0 & 0 & -2 & 0\\0 & 0 & 0 & 1\end{matrix}\right].\]
\step{First solve $Ly=b$ by forward substitution.}
\begin{align*}y_{1}&=1\\
y_{2}&=-2-\bigl[(-2)(1)\bigr]=0\\
y_{3}&=-1-\bigl[(3)(1)+(-3)(0)\bigr]=-4\\
y_{4}&=2-\bigl[(-5)(1)+(4)(0)+(-1)(-4)\bigr]=3\end{align*}\step{Then solve $Ux=y$ by backward substitution.}
\begin{align*}x_{4}&=3=3\\
x_{3}&=\frac{-4}{-2}=2\\
x_{2}&=\frac{0-\bigl[(12)(3)\bigr]}{3}=-12\\
x_{1}&=1-\bigl[(3)(-12)+(2)(2)\bigr]=33\end{align*}\answer{$x=\left[\begin{matrix}33\\-12\\2\\3\end{matrix}\right]$. Multiplying by the original $A$
returns the displayed $b$, an independent check of the solution.}

\exercise{2.5.7}{Record the elimination multipliers to construct $L$}
Start with the given matrix
\[A=\left[\begin{matrix}2 & 5\\-3 & -4\end{matrix}\right].\]
Use the following row replacements, in the order listed. The multiplier
subtracted in each replacement is recorded below the relevant diagonal
entry of $L$:
\begin{align*}R_{2}&\leftarrow R_{2}+(\frac{3}{2})R_{1}\end{align*}The resulting factors are
\[L=\left[\begin{matrix}1 & 0\\- \frac{3}{2} & 1\end{matrix}\right].\]
\[U=\left[\begin{matrix}2 & 5\\0 & \frac{7}{2}\end{matrix}\right].\]Here the diagonal of $L$ consists of ones and all its entries above
the diagonal are zero. The displayed $U$ is in echelon form.
Direct multiplication verifies $LU=A$ exactly.

\exercise{2.5.8}{Record the elimination multipliers to construct $L$}
Start with the given matrix
\[A=\left[\begin{matrix}6 & 4\\12 & 5\end{matrix}\right].\]
Use the following row replacements, in the order listed. The multiplier
subtracted in each replacement is recorded below the relevant diagonal
entry of $L$:
\begin{align*}R_{2}&\leftarrow R_{2}-(2)R_{1}\end{align*}The resulting factors are
\[L=\left[\begin{matrix}1 & 0\\2 & 1\end{matrix}\right].\]
\[U=\left[\begin{matrix}6 & 4\\0 & -3\end{matrix}\right].\]Here the diagonal of $L$ consists of ones and all its entries above
the diagonal are zero. The displayed $U$ is in echelon form.
Direct multiplication verifies $LU=A$ exactly.

\exercise{2.5.9}{Record the elimination multipliers to construct $L$}
Start with the given matrix
\[A=\left[\begin{matrix}3 & 1 & 2\\-9 & 0 & -4\\9 & 9 & 14\end{matrix}\right].\]
Use the following row replacements, in the order listed. The multiplier
subtracted in each replacement is recorded below the relevant diagonal
entry of $L$:
\begin{align*}R_{2}&\leftarrow R_{2}+(3)R_{1}\\
R_{3}&\leftarrow R_{3}-(3)R_{1}\\
R_{3}&\leftarrow R_{3}-(2)R_{2}\end{align*}The resulting factors are
\[L=\left[\begin{matrix}1 & 0 & 0\\-3 & 1 & 0\\3 & 2 & 1\end{matrix}\right].\]
\[U=\left[\begin{matrix}3 & 1 & 2\\0 & 3 & 2\\0 & 0 & 4\end{matrix}\right].\]Here the diagonal of $L$ consists of ones and all its entries above
the diagonal are zero. The displayed $U$ is in echelon form.
Direct multiplication verifies $LU=A$ exactly.

\exercise{2.5.10}{Record the elimination multipliers to construct $L$}
Start with the given matrix
\[A=\left[\begin{matrix}-5 & 0 & 4\\10 & 2 & -5\\10 & 10 & 16\end{matrix}\right].\]
Use the following row replacements, in the order listed. The multiplier
subtracted in each replacement is recorded below the relevant diagonal
entry of $L$:
\begin{align*}R_{2}&\leftarrow R_{2}+(2)R_{1}\\
R_{3}&\leftarrow R_{3}+(2)R_{1}\\
R_{3}&\leftarrow R_{3}-(5)R_{2}\end{align*}The resulting factors are
\[L=\left[\begin{matrix}1 & 0 & 0\\-2 & 1 & 0\\-2 & 5 & 1\end{matrix}\right].\]
\[U=\left[\begin{matrix}-5 & 0 & 4\\0 & 2 & 3\\0 & 0 & 9\end{matrix}\right].\]Here the diagonal of $L$ consists of ones and all its entries above
the diagonal are zero. The displayed $U$ is in echelon form.
Direct multiplication verifies $LU=A$ exactly.

\exercise{2.5.11}{Record the elimination multipliers to construct $L$}
Start with the given matrix
\[A=\left[\begin{matrix}3 & 7 & 2\\6 & 19 & 4\\-3 & -2 & 3\end{matrix}\right].\]
Use the following row replacements, in the order listed. The multiplier
subtracted in each replacement is recorded below the relevant diagonal
entry of $L$:
\begin{align*}R_{2}&\leftarrow R_{2}-(2)R_{1}\\
R_{3}&\leftarrow R_{3}+R_{1}\\
R_{3}&\leftarrow R_{3}-R_{2}\end{align*}The resulting factors are
\[L=\left[\begin{matrix}1 & 0 & 0\\2 & 1 & 0\\-1 & 1 & 1\end{matrix}\right].\]
\[U=\left[\begin{matrix}3 & 7 & 2\\0 & 5 & 0\\0 & 0 & 5\end{matrix}\right].\]Here the diagonal of $L$ consists of ones and all its entries above
the diagonal are zero. The displayed $U$ is in echelon form.
Direct multiplication verifies $LU=A$ exactly.

\exercise{2.5.12}{Record the elimination multipliers to construct $L$}
Start with the given matrix
\[A=\left[\begin{matrix}2 & 3 & 2\\4 & 13 & 9\\-6 & 5 & 4\end{matrix}\right].\]
Use the following row replacements, in the order listed. The multiplier
subtracted in each replacement is recorded below the relevant diagonal
entry of $L$:
\begin{align*}R_{2}&\leftarrow R_{2}-(2)R_{1}\\
R_{3}&\leftarrow R_{3}+(3)R_{1}\\
R_{3}&\leftarrow R_{3}-(2)R_{2}\end{align*}The resulting factors are
\[L=\left[\begin{matrix}1 & 0 & 0\\2 & 1 & 0\\-3 & 2 & 1\end{matrix}\right].\]
\[U=\left[\begin{matrix}2 & 3 & 2\\0 & 7 & 5\\0 & 0 & 0\end{matrix}\right].\]Here the diagonal of $L$ consists of ones and all its entries above
the diagonal are zero. The displayed $U$ is in echelon form.
Direct multiplication verifies $LU=A$ exactly. A zero final row does not
invalidate an $LU$ factorization; the matrix need not be invertible.

\exercise{2.5.13}{Record the elimination multipliers to construct $L$}
Start with the given matrix
\[A=\left[\begin{matrix}1 & 3 & -5 & -3\\-1 & -5 & 8 & 4\\4 & 2 & -5 & -7\\-2 & -4 & 7 & 5\end{matrix}\right].\]
Use the following row replacements, in the order listed. The multiplier
subtracted in each replacement is recorded below the relevant diagonal
entry of $L$:
\begin{align*}R_{2}&\leftarrow R_{2}+R_{1}\\
R_{3}&\leftarrow R_{3}-(4)R_{1}\\
R_{4}&\leftarrow R_{4}+(2)R_{1}\\
R_{3}&\leftarrow R_{3}-(5)R_{2}\\
R_{4}&\leftarrow R_{4}+R_{2}\end{align*}The resulting factors are
\[L=\left[\begin{matrix}1 & 0 & 0 & 0\\-1 & 1 & 0 & 0\\4 & 5 & 1 & 0\\-2 & -1 & 0 & 1\end{matrix}\right].\]
\[U=\left[\begin{matrix}1 & 3 & -5 & -3\\0 & -2 & 3 & 1\\0 & 0 & 0 & 0\\0 & 0 & 0 & 0\end{matrix}\right].\]Here the diagonal of $L$ consists of ones and all its entries above
the diagonal are zero. The displayed $U$ is in echelon form.
Direct multiplication verifies $LU=A$ exactly. A zero final row does not
invalidate an $LU$ factorization; the matrix need not be invertible.

\exercise{2.5.14}{Record the elimination multipliers to construct $L$}
Start with the given matrix
\[A=\left[\begin{matrix}1 & 3 & 1 & 5\\5 & 20 & 6 & 31\\-2 & -1 & -1 & -4\\-1 & 7 & 1 & 7\end{matrix}\right].\]
Use the following row replacements, in the order listed. The multiplier
subtracted in each replacement is recorded below the relevant diagonal
entry of $L$:
\begin{align*}R_{2}&\leftarrow R_{2}-(5)R_{1}\\
R_{3}&\leftarrow R_{3}+(2)R_{1}\\
R_{4}&\leftarrow R_{4}+R_{1}\\
R_{3}&\leftarrow R_{3}-R_{2}\\
R_{4}&\leftarrow R_{4}-(2)R_{2}\end{align*}The resulting factors are
\[L=\left[\begin{matrix}1 & 0 & 0 & 0\\5 & 1 & 0 & 0\\-2 & 1 & 1 & 0\\-1 & 2 & 0 & 1\end{matrix}\right].\]
\[U=\left[\begin{matrix}1 & 3 & 1 & 5\\0 & 5 & 1 & 6\\0 & 0 & 0 & 0\\0 & 0 & 0 & 0\end{matrix}\right].\]Here the diagonal of $L$ consists of ones and all its entries above
the diagonal are zero. The displayed $U$ is in echelon form.
Direct multiplication verifies $LU=A$ exactly. A zero final row does not
invalidate an $LU$ factorization; the matrix need not be invertible.

\exercise{2.5.15}{Record the elimination multipliers to construct $L$}
Start with the given matrix
\[A=\left[\begin{matrix}2 & 0 & 5 & 2\\-6 & 3 & -13 & -3\\4 & 9 & 16 & 17\end{matrix}\right].\]
Use the following row replacements, in the order listed. The multiplier
subtracted in each replacement is recorded below the relevant diagonal
entry of $L$:
\begin{align*}R_{2}&\leftarrow R_{2}+(3)R_{1}\\
R_{3}&\leftarrow R_{3}-(2)R_{1}\\
R_{3}&\leftarrow R_{3}-(3)R_{2}\end{align*}The resulting factors are
\[L=\left[\begin{matrix}1 & 0 & 0\\-3 & 1 & 0\\2 & 3 & 1\end{matrix}\right].\]
\[U=\left[\begin{matrix}2 & 0 & 5 & 2\\0 & 3 & 2 & 3\\0 & 0 & 0 & 4\end{matrix}\right].\]Here the diagonal of $L$ consists of ones and all its entries above
the diagonal are zero. The displayed $U$ is in echelon form.
Direct multiplication verifies $LU=A$ exactly. Because $A$ is rectangular, $L$ is
$3\times3$ and $U$ has the same
$3\times4$ size as $A$.

\exercise{2.5.16}{Record the elimination multipliers to construct $L$}
Start with the given matrix
\[A=\left[\begin{matrix}2 & -3 & 4\\-4 & 8 & -7\\6 & -5 & 14\\-6 & 9 & -12\\8 & -6 & 19\end{matrix}\right].\]
Use the following row replacements, in the order listed. The multiplier
subtracted in each replacement is recorded below the relevant diagonal
entry of $L$:
\begin{align*}R_{2}&\leftarrow R_{2}+(2)R_{1}\\
R_{3}&\leftarrow R_{3}-(3)R_{1}\\
R_{4}&\leftarrow R_{4}+(3)R_{1}\\
R_{5}&\leftarrow R_{5}-(4)R_{1}\\
R_{3}&\leftarrow R_{3}-(2)R_{2}\\
R_{5}&\leftarrow R_{5}-(3)R_{2}\end{align*}The resulting factors are
\[L=\left[\begin{matrix}1 & 0 & 0 & 0 & 0\\-2 & 1 & 0 & 0 & 0\\3 & 2 & 1 & 0 & 0\\-3 & 0 & 0 & 1 & 0\\4 & 3 & 0 & 0 & 1\end{matrix}\right].\]
\[U=\left[\begin{matrix}2 & -3 & 4\\0 & 2 & 1\\0 & 0 & 0\\0 & 0 & 0\\0 & 0 & 0\end{matrix}\right].\]Here the diagonal of $L$ consists of ones and all its entries above
the diagonal are zero. The displayed $U$ is in echelon form.
Direct multiplication verifies $LU=A$ exactly. A zero final row does not
invalidate an $LU$ factorization; the matrix need not be invertible. Because $A$ is rectangular, $L$ is
$5\times5$ and $U$ has the same
$5\times3$ size as $A$.

\exercise{2.5.17}{Invert the two triangular factors in reverse order}
Use the factors from Exercise 2.5.2. Solving $Lz=e_j$
column by column (forward substitution), and $Uw=e_j$ column by
column (backward substitution), gives
\[L^{-1}=\left[\begin{matrix}1 & 0 & 0\\2 & 1 & 0\\-2 & -1 & 1\end{matrix}\right],\qquad U^{-1}=\left[\begin{matrix}\frac{1}{2} & - \frac{3}{4} & -2\\0 & - \frac{1}{4} & -1\\0 & 0 & - \frac{1}{2}\end{matrix}\right].\]
This is the augmented-inverse algorithm: reduce $[L\mid I]$ and
$[U\mid I]$ to $[I\mid L^{-1}]$ and $[I\mid U^{-1}]$ using their
triangular equations. Each column is checked by multiplying it by its
original factor and obtaining the corresponding coordinate vector.
Since $(LU)^{-1}=U^{-1}L^{-1}$, the required inverse is
\[A^{-1}=U^{-1}L^{-1}=\left[\begin{matrix}3 & \frac{5}{4} & -2\\\frac{3}{2} & \frac{3}{4} & -1\\1 & \frac{1}{2} & - \frac{1}{2}\end{matrix}\right].\]
Both $AA^{-1}$ and $A^{-1}A$ equal the identity. The order of the
two inverse factors matters.

\exercise{2.5.18}{Invert the two triangular factors in reverse order}
Use the factors from Exercise 2.5.3. Solving $Lz=e_j$
column by column (forward substitution), and $Uw=e_j$ column by
column (backward substitution), gives
\[L^{-1}=\left[\begin{matrix}1 & 0 & 0\\2 & 1 & 0\\-1 & 1 & 1\end{matrix}\right],\qquad U^{-1}=\left[\begin{matrix}\frac{1}{2} & - \frac{2}{3} & 3\\0 & - \frac{1}{3} & 2\\0 & 0 & 1\end{matrix}\right].\]
This is the augmented-inverse algorithm: reduce $[L\mid I]$ and
$[U\mid I]$ to $[I\mid L^{-1}]$ and $[I\mid U^{-1}]$ using their
triangular equations. Each column is checked by multiplying it by its
original factor and obtaining the corresponding coordinate vector.
Since $(LU)^{-1}=U^{-1}L^{-1}$, the required inverse is
\[A^{-1}=U^{-1}L^{-1}=\left[\begin{matrix}- \frac{23}{6} & \frac{7}{3} & 3\\- \frac{8}{3} & \frac{5}{3} & 2\\-1 & 1 & 1\end{matrix}\right].\]
Both $AA^{-1}$ and $A^{-1}A$ equal the identity. The order of the
two inverse factors matters.

\exercise{2.5.19}{Invert a triangular matrix without losing triangular form}
Every diagonal entry is nonzero, so forward substitution solves
$Ax=b$ uniquely for every $b$. Hence $A$ is invertible by the
Invertible Matrix Theorem.
To follow the suggested row-operation argument, clear entries below
the successive diagonal pivots using replacements by earlier rows,
then scale the diagonal pivots to $1$. No swaps or additions of later
rows into earlier rows are needed. Applying these same operations to
$I$ keeps it lower triangular: an earlier lower triangular row has
zeros in every column beyond its own index, hence also beyond the
index of the later row it is added to. Scaling preserves these zeros.
Thus reduction of $[A\mid I]$ to $[I\mid A^{-1}]$ leaves the right
block lower triangular. This proves the requested structure as well
as invertibility.

\exercise{2.5.20}{Reduce the unit lower factor to the identity}
Because $L$ is unit lower triangular, its diagonal entries are already
$1$. Clear below each diagonal using only replacements of a later row
by itself minus a multiple of an earlier row. Let $E$ be the product
of these elementary row-operation matrices. Then $EL=I$.
Applying the same replacements to $A=LU$ gives
\[EA=ELU=IU=U.\]
Thus only row replacements are needed, even if $A$ or $U$ is
rectangular or singular.

\exercise{2.5.21}{One sequence of row operations acts on every factorization}
Let $E$ be the product of all elementary matrices for the stated
operations. Since they reduce $B$ to $I$, $EB=I$; therefore
\[EA=E(BC)=(EB)C=IC=C.\]
This proves the claim for any such sequence, not merely a particular
algorithm. The converse can fail: take $A=0$, any invertible $B\ne I$,
and $C=0$. The empty sequence already leaves $A=C$, but does not
change $B$ into $I$.

\exercise{2.5.22}{Discard the zero rows and their unused multiplier columns}
For the Practice Problem factors, take the first three columns of
$L$ and the first three rows of $U$:
\[B=\left[\begin{matrix}1 & 0 & 0\\3 & 1 & 0\\1 & -1 & 1\\2 & 2 & -1\\-3 & -3 & 2\end{matrix}\right],\qquad C=\left[\begin{matrix}2 & -4 & -2 & 3\\0 & 3 & 1 & -1\\0 & 0 & 0 & 5\end{matrix}\right].\]
They have the required sizes $5\times3$ and $3\times4$, and their
product is $A$. To see why, write
$L=[L_1\ L_2]$ and $U=\begin{bmatrix}U_1\\0\end{bmatrix}$.
Block multiplication gives $LU=L_1U_1+L_2(0)=L_1U_1$.
The same argument works for any $m\times n$ matrix whose $U$ has
exactly three nonzero rows: use an $m\times3$ factor and a
$3\times n$ factor. Those rows are first because $U$ is echelon.

\exercise{2.5.23}{A compact representation by four outer products}
\begin{enumerate}[label=(\alph*)]
\item Write $C=[c_1\ c_2\ c_3\ c_4]$ and let the rows of $D$ be
$d_1^T,\ldots,d_4^T$. The column--row rule gives
$A=c_1d_1^T+c_2d_2^T+c_3d_3^T+c_4d_4^T$.
Each term is an $m\times n$ outer product.
\item A $400\times100$ array stores $40{,}000$ numbers.
The factors require $400(4)+4(100)=2{,}000$ numbers, a factor
of $20$ less storage, ignoring small bookkeeping costs.
For a product $Ax$, first compute $Dx$ (four entries) and then
$C(Dx)$; storing the full matrix is unnecessary for that use.
\end{enumerate}

\exercise{2.5.24}{Solve a QR-factored system by an orthogonal multiplication}
From $Q^TQ=I$ and the fact that $Q$ is square, $Q$ is invertible
with inverse $Q^T$. Thus $QRx=b$ is equivalent to
$Rx=Q^Tb$. Compute $c=Q^Tb$, then solve the upper triangular
system $Rx=c$ by backward substitution. Invertibility of $R$
guarantees a nonzero diagonal and a unique solution.
\answer{For every $b$, this process yields the unique $x$. It requires
no explicit inverse of $R$.}

\exercise{2.5.25}{Invert a singular-value factorization with positive diagonal}
The square matrices $U,V$ are invertible because $U^TU=V^TV=I$;
their inverses are $U^T,V^T$. Each $\sigma_j>0$, so
$D^{-1}=\operatorname{diag}(1/\sigma_1,\ldots,1/\sigma_n)$ exists.
Reverse the factor order:
\[A^{-1}=(V^T)^{-1}D^{-1}U^{-1}=VD^{-1}U^T.\]
The product check is
$AA^{-1}=UDV^TVD^{-1}U^T=UU^T=I$, and the reverse order
similarly gives $VV^T=I$. Strict positivity is stronger than the
nonzero condition needed for the inverse, but supplies it here.

\exercise{2.5.26}{Cancel neighboring change-of-basis factors in powers}
Associativity and $P^{-1}P=I$ give
\[A^2=(PDP^{-1})(PDP^{-1})=PD^2P^{-1},\]
\[A^3=PD^3P^{-1},\qquad A^k=PD^kP^{-1}.\]
The general formula follows inductively by multiplying the formula
for $A^k$ by $PDP^{-1}$. Since $D=\operatorname{diag}(2,3,1)$,
\[A^2=P\operatorname{diag}(4,9,1)P^{-1},\quad
A^3=P\operatorname{diag}(8,27,1)P^{-1},\]
\[A^k=P\operatorname{diag}(2^k,3^k,1)P^{-1}.\]
Only scalar powers of three diagonal entries need repeated powering.

\textbf{Circuit convention for the next four exercises.}
A series resistor of resistance $R$ has transfer matrix
$S(R)=\begin{bmatrix}1&-R\\0&1\end{bmatrix}$; a shunt resistor has
$T(R)=\begin{bmatrix}1&0\\-1/R&1\end{bmatrix}$.
The state is $(v,i)^T$ and is mapped from input to output.
If the signal passes through circuit 1 and then circuit 2, the total
matrix is $A_2A_1$, so physical order is reversed in the matrix product.

\exercise{2.5.27}{Two different networks with the specified input and output}
\step{First design: a series resistor followed by a shunt.}
A series resistance $R=1/2$ changes $(12,6)^T$ to
$(12-(1/2)6,6)^T=(9,6)^T$. A shunt resistance $9/2$ then
changes this to $(9,6-9/(9/2))^T=(9,4)^T$.
Its matrix is $T(9/2)S(1/2)$.
\step{Second design: a shunt followed by a series resistor.}
A shunt resistance $6$ changes the input to $(12,6-12/6)^T=(12,4)^T$.
A series resistance $3/4$ then gives $(12-(3/4)4,4)^T=(9,4)^T$.
Its matrix is $S(3/4)T(6)$. All four resistances are positive and
both calculations verify the required output, despite the different
physical orders and resistance values.

\exercise{2.5.28}{Add conductances for three successive shunt circuits}
Multiplying two shunt matrices adds their lower-left entries.
Thus three of them give
\[T(R_3)T(R_2)T(R_1)
=\begin{bmatrix}1&0\\-(1/R_1+1/R_2+1/R_3)&1\end{bmatrix}.\]
This is $T(R)$ with
\[R=\frac1{1/R_1+1/R_2+1/R_3}.
\]
The positive-resistance assumption makes the denominator positive.
One may also write $R=R_1R_2R_3/(R_1R_2+R_1R_3+R_2R_3)$.

\exercise{2.5.29}{Factor the given series--shunt--series network}
\step{(a) Multiply in the physical order's reverse.}
The signal passes through $R_1$ in series, $R_2$ as a shunt,
and $R_3$ in series. Hence
\[S(R_3)T(R_2)S(R_1)
=\begin{bmatrix}
1+R_3/R_2&-R_1-R_3-R_1R_3/R_2\\
-1/R_2&1+R_1/R_2
\end{bmatrix}.\]
For example, multiplying the right two factors first gives
$\begin{bmatrix}1&-R_1\\-1/R_2&1+R_1/R_2\end{bmatrix}$;
the final left multiplication gives the displayed first row.
\step{(b) Match entries of the target matrix.}
The lower-left entry $-1/3$ forces $R_2=3$.
The lower-right entry $5/3$ then gives $R_1=2$.
The upper-left entry $3$ gives $R_3=6$.
The remaining entry checks: $-2-6-2(6)/3=-12$.
\answer{Use a $2$-ohm series resistor, a $3$-ohm shunt, then a
$6$-ohm series resistor. Their product is exactly the specified matrix.}

\exercise{2.5.30}{A shunt--series--shunt realization of the same matrix}
Try a shunt of resistance $a$, then a series resistance $b$, then
a shunt of resistance $c$. Multiplying gives
\[T(c)S(b)T(a)=\begin{bmatrix}
1+b/a&-b\\-1/a-(1+b/a)/c&1+b/c
\end{bmatrix}.\]
Match the target's upper-right entry to obtain $b=12$.
Then $1+12/a=3$ gives $a=6$, and $1+12/c=5/3$ gives $c=18$.
The lower-left entry checks as $-1/6-3/18=-1/3$.
\answer{The distinct physical design is a $6$-ohm shunt, a
$12$-ohm series resistor, and an $18$-ohm shunt, with exactly the
same overall transfer matrix as in Exercise 2.5.29.}

\exercise{2.5.31}{Factor and solve the supplied eight-node heat plate}
All omitted entries in the printed matrix are zero; the actual
coefficient matrix used here is
\[A=\left[\begin{matrix}4 & -1 & -1 & 0 & 0 & 0 & 0 & 0\\-1 & 4 & 0 & -1 & 0 & 0 & 0 & 0\\-1 & 0 & 4 & -1 & -1 & 0 & 0 & 0\\0 & -1 & -1 & 4 & 0 & -1 & 0 & 0\\0 & 0 & -1 & 0 & 4 & -1 & -1 & 0\\0 & 0 & 0 & -1 & -1 & 4 & 0 & -1\\0 & 0 & 0 & 0 & -1 & 0 & 4 & -1\\0 & 0 & 0 & 0 & 0 & -1 & -1 & 4\end{matrix}\right].\]
\step{(a) Perform elimination without row swaps.}
The successive pivots are the diagonal entries of the following $U$.
At pivot row $j$, subtract $l_{ij}$ times that row from row $i$;
the multipliers $l_{ij}$ are precisely the displayed entries of $L$.
\[L=\left[\begin{matrix}1 & 0 & 0 & 0 & 0 & 0 & 0 & 0\\- \frac{1}{4} & 1 & 0 & 0 & 0 & 0 & 0 & 0\\- \frac{1}{4} & - \frac{1}{15} & 1 & 0 & 0 & 0 & 0 & 0\\0 & - \frac{4}{15} & - \frac{2}{7} & 1 & 0 & 0 & 0 & 0\\0 & 0 & - \frac{15}{56} & - \frac{1}{12} & 1 & 0 & 0 & 0\\0 & 0 & 0 & - \frac{7}{24} & - \frac{26}{89} & 1 & 0 & 0\\0 & 0 & 0 & 0 & - \frac{24}{89} & - \frac{208}{2415} & 1 & 0\\0 & 0 & 0 & 0 & 0 & - \frac{712}{2415} & - \frac{2623}{8948} & 1\end{matrix}\right].\]
\[U=\left[\begin{matrix}4 & -1 & -1 & 0 & 0 & 0 & 0 & 0\\0 & \frac{15}{4} & - \frac{1}{4} & -1 & 0 & 0 & 0 & 0\\0 & 0 & \frac{56}{15} & - \frac{16}{15} & -1 & 0 & 0 & 0\\0 & 0 & 0 & \frac{24}{7} & - \frac{2}{7} & -1 & 0 & 0\\0 & 0 & 0 & 0 & \frac{89}{24} & - \frac{13}{12} & -1 & 0\\0 & 0 & 0 & 0 & 0 & \frac{2415}{712} & - \frac{26}{89} & -1\\0 & 0 & 0 & 0 & 0 & 0 & \frac{8948}{2415} & - \frac{2623}{2415}\\0 & 0 & 0 & 0 & 0 & 0 & 0 & \frac{30305}{8948}\end{matrix}\right].\]
Both factors have zeros more than two positions below or above their
main diagonals. Exact multiplication gives $LU-A=0_{8\times8}$,
checking every entry rather than just the diagonals.
\step{(b) Solve using the factors.}
For $b=(5,15,0,10,0,10,20,30)^T$, forward substitution gives
\[y=\left[\begin{matrix}5\\\frac{65}{4}\\\frac{7}{3}\\15\\\frac{15}{8}\\\frac{10625}{712}\\\frac{10525}{483}\\\frac{364965}{8948}\end{matrix}\right].\]
The rule is $y_i=b_i-\sum_{j<i}l_{ij}y_j$, starting at row 1.
Backward substitution uses
$x_i=(y_i-\sum_{j>i}u_{ij}x_j)/u_{ii}$, starting at row 8.
It gives
\[x=\frac1{209}\begin{bmatrix}827\\1377\\886\\1546\\1171\\1831\\1967\\2517\end{bmatrix}.
\]
Substitution verifies $Ax=b$ exactly. For instance row 1 gives
$(4(827)-1377-886)/209=1045/209=5$.
\step{(c) Compare the inverse with the banded factors.}
The exact inverse is
\[A^{-1}=\frac1{30305}\left[\begin{matrix}8948 & 2623 & 2864 & 1544 & 964 & 689 & 303 & 248\\2623 & 8948 & 1544 & 2864 & 689 & 964 & 248 & 303\\2864 & 1544 & 9912 & 3312 & 3167 & 1792 & 964 & 689\\1544 & 2864 & 3312 & 9912 & 1792 & 3167 & 689 & 964\\964 & 689 & 3167 & 1792 & 9912 & 3312 & 2864 & 1544\\689 & 964 & 1792 & 3167 & 3312 & 9912 & 1544 & 2864\\303 & 248 & 964 & 689 & 2864 & 1544 & 8948 & 2623\\248 & 303 & 689 & 964 & 1544 & 2864 & 2623 & 8948\end{matrix}\right].\]
Every entry is nonzero. Unlike $L,U$, it has no empty outer bands.
For an $n\times n$ matrix with fixed band width, storing its factors
requires a number of entries proportional to $n$, whereas a dense
inverse requires $n^2$. This example verifies the stated storage
advantage with actual factors and an actual inverse.

\exercise{2.5.32}{Reuse one bidiagonal factorization for four time steps}
\step{(a) Set $C=1$ before factoring.}
Then $A$ has diagonal $3$ and adjacent off-diagonal entries $-1$.
Its factors are
\[L=\left[\begin{matrix}1 & 0 & 0 & 0\\- \frac{1}{3} & 1 & 0 & 0\\0 & - \frac{3}{8} & 1 & 0\\0 & 0 & - \frac{8}{21} & 1\end{matrix}\right].\]
\[U=\left[\begin{matrix}3 & -1 & 0 & 0\\0 & \frac{8}{3} & -1 & 0\\0 & 0 & \frac{21}{8} & -1\\0 & 0 & 0 & \frac{55}{21}\end{matrix}\right].\]
Their off-diagonal nonzero entries lie on just one neighboring diagonal,
so both factors are bidiagonal. Multiplication checks $LU=A$.
\step{(b) Solve the same two triangular systems at each step.}
Starting with $t_0=(10,15,15,10)^T$, repeatedly solve
$Ly=t_k$ and $Ut_{k+1}=y$. The results are
\[\begin{array}{c|rrrr}
k&(t_k)_1&(t_k)_2&(t_k)_3&(t_k)_4\\\hline
1&7&11&11&7\\2&5&8&8&5\\
3&18/5&29/5&29/5&18/5\\4&13/5&21/5&21/5&13/5
\end{array}.\]
A short independent check uses symmetry. For a vector $(a,b,b,a)^T$,
$A(a,b,b,a)^T=(3a-b,2b-a,2b-a,3a-b)^T$.
Thus if $(a_0,b_0)$ is the preceding pair, the next pair is
$a=(2a_0+b_0)/5$, $b=(a_0+3b_0)/5$.
These formulas reproduce all four rows and verify
$At_{k+1}=t_k$ at each step. The factorization is computed once,
not recomputed for each new temperature distribution.
